\documentclass[11pt,a4paper,landscape]{article}

\usepackage[utf8]{inputenc}
\usepackage[T1]{fontenc}
\usepackage[french]{babel}

\usepackage{styles/couleurs}
\usepackage{styles/boites}
\usepackage{styles/maths}
\usepackage{styles/figures}

\usepackage{amsmath}
\usepackage{amssymb}
\usepackage[most]{tcolorbox}
\usepackage{geometry}

\geometry{
  a4paper,
  landscape,
  left=9mm,
  right=9mm,
  top=8mm,
  bottom=8mm
}

\pagestyle{empty}
\setlength{\parindent}{0pt}
\setlength{\parskip}{0pt}

\begin{document}

% =========================================================
% TITRE
% =========================================================

\begin{center}
{\sffamily\LARGE\bfseries\color{TitreC}
AP -- Je m'entraîne à calculer avec Pythagore}

\vspace{0.1cm}

{\sffamily\large\color{SavoirC}
Correction}
\end{center}

\vspace{0.3cm}

% =========================================================
% LIGNE 1
% =========================================================

\noindent
\begin{minipage}[t]{0.235\textwidth}

\begin{tcolorbox}[
  colback=white,
  colframe=ActionC!50,
  boxrule=0.6pt,
  arc=1.5mm,
  height=7cm
]

\begin{center}
{\sffamily\bfseries\color{ActionC}Triangle ABC}
\end{center}

\[
8\times8=64\text{ cm}^2
\]

\[
15\times15=225\text{ cm}^2
\]

\[
64+225=289\text{ cm}^2
\]

\[
17\times17=289\text{ cm}^2
\]

\[
\boxed{BC=17\text{ cm}}
\]

\end{tcolorbox}

\end{minipage}
\hfill
%
\begin{minipage}[t]{0.235\textwidth}

\begin{tcolorbox}[
  colback=white,
  colframe=ActionC!50,
  boxrule=0.6pt,
  arc=1.5mm,
  height=7cm
]

\begin{center}
{\sffamily\bfseries\color{ActionC}Triangle DEF}
\end{center}

\[
41\times41=1681\text{ cm}^2
\]

\[
9\times9=81\text{ cm}^2
\]

\[
1681-81=1600\text{ cm}^2
\]

\[
40\times40=1600\text{ cm}^2
\]

\[
\boxed{EF=40\text{ cm}}
\]

\end{tcolorbox}

\end{minipage}
\hfill
%
\begin{minipage}[t]{0.235\textwidth}

\begin{tcolorbox}[
  colback=white,
  colframe=ActionC!50,
  boxrule=0.6pt,
  arc=1.5mm,
  height=7cm
]

\begin{center}
{\sffamily\bfseries\color{ActionC}Triangle GHI}
\end{center}

\[
12\times12=144\text{ cm}^2
\]

\[
35\times35=1225\text{ cm}^2
\]

\[
144+1225=1369\text{ cm}^2
\]

\[
37\times37=1369\text{ cm}^2
\]

\[
\boxed{HI=37\text{ cm}}
\]

\end{tcolorbox}

\end{minipage}
\hfill
%
\begin{minipage}[t]{0.235\textwidth}

\begin{tcolorbox}[
  colback=white,
  colframe=ActionC!50,
  boxrule=0.6pt,
  arc=1.5mm,
  height=7cm
]

\begin{center}
{\sffamily\bfseries\color{ActionC}Triangle JKL}
\end{center}

\[
65\times65=4225\text{ cm}^2
\]

\[
33\times33=1089\text{ cm}^2
\]

\[
4225-1089=3136\text{ cm}^2
\]

\[
56\times56=3136\text{ cm}^2
\]

\[
\boxed{JL=56\text{ cm}}
\]

\end{tcolorbox}

\end{minipage}

\vspace{0.35cm}

% =========================================================
% LIGNE 2
% =========================================================

\noindent
\begin{minipage}[t]{0.235\textwidth}

\begin{tcolorbox}[
  colback=white,
  colframe=SavoirC!50,
  boxrule=0.6pt,
  arc=1.5mm,
  height=7cm
]

\begin{center}
{\sffamily\bfseries\color{SavoirC}Triangle MNP}
\end{center}

\[
2{,}4\times2{,}4=5{,}76\text{ cm}^2
\]

\[
4{,}5\times4{,}5=20{,}25\text{ cm}^2
\]

\[
5{,}76+20{,}25=26{,}01\text{ cm}^2
\]

\[
5{,}1\times5{,}1=26{,}01\text{ cm}^2
\]

\[
\boxed{NP=5{,}1\text{ cm}}
\]

\end{tcolorbox}

\end{minipage}
\hfill
%
\begin{minipage}[t]{0.235\textwidth}

\begin{tcolorbox}[
  colback=white,
  colframe=SavoirC!50,
  boxrule=0.6pt,
  arc=1.5mm,
  height=7cm
]

\begin{center}
{\sffamily\bfseries\color{SavoirC}Triangle QRS}
\end{center}

\[
1{,}8\times1{,}8=3{,}24\text{ cm}^2
\]

\[
8\times8=64\text{ cm}^2
\]

\[
3{,}24+64=67{,}24\text{ cm}^2
\]

\[
8{,}2\times8{,}2=67{,}24\text{ cm}^2
\]

\[
\boxed{RS=8{,}2\text{ cm}}
\]

\end{tcolorbox}

\end{minipage}
\hfill
%
\begin{minipage}[t]{0.235\textwidth}

\begin{tcolorbox}[
  colback=white,
  colframe=SavoirC!50,
  boxrule=0.6pt,
  arc=1.5mm,
  height=7cm
]

\begin{center}
{\sffamily\bfseries\color{SavoirC}Triangle TUV}
\end{center}

\[
39\times39=1521\text{ cm}^2
\]

\[
80\times80=6400\text{ cm}^2
\]

\[
1521+6400=7921\text{ cm}^2
\]

\[
89\times89=7921\text{ cm}^2
\]

\[
\boxed{TV=89\text{ cm}}
\]

\end{tcolorbox}

\end{minipage}
\hfill
%
\begin{minipage}[t]{0.235\textwidth}

\begin{tcolorbox}[
  colback=SavoirBg,
  colframe=SavoirC!65,
  boxrule=0.8pt,
  arc=1.5mm,
  height=7cm
]

\begin{center}
{\sffamily\bfseries\color{SavoirC}Triangle EFG}
\end{center}

\[
5\times5=25\text{ cm}^2
\]

\[
5\times5=25\text{ cm}^2
\]

\[
25+25=50\text{ cm}^2
\]

\[
7\times7=49\text{ cm}^2
\qquad
8\times8=64\text{ cm}^2
\]

\[
\boxed{7\text{ cm}<EG<8\text{ cm}}
\]

\begin{center}
\small
Quel nombre positif multiplié par lui-même donne \(50\) ?
\end{center}

\end{tcolorbox}

\end{minipage}

\end{document}